% ============================================================================
%  PRECÁLCULO — Sesión 13
%  Semana 13: 2 al 7 de noviembre de 2026   (lunes 2: festivo)
%  3.1 Operaciones algebraicas · 3.2 Productos notables
%  Duración prevista: 3 horas
% ============================================================================
\documentclass[aspectratio=169,10pt]{beamer}

\usepackage[utf8]{inputenc}
\IfFileExists{spanish.ldf}{\usepackage[spanish,es-noshorthands]{babel}}{}
\usepackage{amsmath,amssymb}
\usepackage{booktabs}
\usepackage{tikz}

\usetheme{Uniminuto}

\uniasignatura{Precálculo}
\uniprofesor{Andrés Felipe Puerta Velez}
\unicorreo{andres.puerta-v@uniminuto.edu.co}
\uninrc{MATE1161}
\unisemana{Sesión 13 — Semana 13 · 2 al 7 de noviembre de 2026}

\newcommand{\listatiny}{\setbeamerfont{itemize/enumerate body}{size=\tiny}\tiny}

\begin{document}

\portada

\begin{frame}{Orden del día}
  \begin{enumerate}
    \item Retroalimentación del Examen 2
    \item \azul{Tema 3.1: lenguaje algebraico y operaciones con polinomios}
    \item Suma, resta, producto y división de polinomios
    \item \azul{Tema 3.2: productos notables}
    \item Ejercicios guiados, taller y cierre
  \end{enumerate}

  \vspace{0.25cm}
  {\footnotesize\color{uniGris} Sesión de 3 horas. Empieza la Unidad 3.
  \destacar{No hay reglas nuevas}: son las mismas propiedades de siempre, pero
  ahora con letras en lugar de números.}
\end{frame}

\begin{frame}{De los números a las letras}
  \scriptsize
  \begin{columns}[T,onlytextwidth]
    \begin{column}{0.48\textwidth}
      \begin{block}{Vocabulario}
        \tiny
        \textbf{Término:} $-5x^{3}y$. Su \azul{coeficiente} es $-5$; su
        \azul{parte literal}, $x^{3}y$; su \azul{grado}, $3+1=4$.

        \vspace{0.06cm}
        \textbf{Términos semejantes:} misma parte literal. \\
        $3x^{2}$ y $-7x^{2}$ lo son; $3x^{2}$ y $3x^{3}$ no.

        \vspace{0.06cm}
        \textbf{Polinomio:} suma de términos. Monomio, binomio, trinomio.
      \end{block}
    \end{column}
    \begin{column}{0.48\textwidth}
      \begin{alertblock}{Por qué se usa el álgebra}
        \tiny
        Porque permite escribir una regla \destacar{una sola vez} y aplicarla a
        infinitos casos.

        \vspace{0.06cm}
        ``El área de un rectángulo cuyo largo excede en 3 al ancho'' es
        $A = x(x+3)$: \azul{una fórmula, no una cuenta}.

        \vspace{0.06cm}
        Es exactamente lo que hace una \destacar{función} en programación.
      \end{alertblock}
    \end{column}
  \end{columns}

  \vspace{0.1cm}
  \begin{exampleblock}{La regla que gobierna todo}
    \scriptsize
    \destacar{Solo se suman términos semejantes.} $3x^{2} + 5x$ no se puede
    reducir, igual que $\sqrt2 + \sqrt3$ no se podía en la sesión 8. En
    cambio, \azul{multiplicar} sí se puede siempre.
  \end{exampleblock}
\end{frame}

\seccion{3.1 Operaciones algebraicas}

\begin{frame}{Suma, resta y producto}
  \scriptsize
  \begin{columns}[T,onlytextwidth]
    \begin{column}{0.50\textwidth}
      \begin{block}{Suma y resta}
        \tiny
        Se reducen términos semejantes:
        \[ (3x^{2}-5x+2) + (x^{2}+7x-9) \]
        \[ = \destacar{4x^{2}+2x-7} \]
        Al restar, \azul{se cambia el signo de todo} el segundo paréntesis:
        \[ (3x^{2}-5x) - (x^{2}-7x) = 2x^{2}+2x \]
      \end{block}
    \end{column}
    \begin{column}{0.46\textwidth}
      \begin{block}{Producto de monomios}
        \tiny
        Coeficientes por coeficientes; letras por letras, \azul{sumando
        exponentes}:
        \[ (4a^{2}b)(3ab^{3}) = \destacar{12a^{3}b^{4}} \]
      \end{block}
    \end{column}
  \end{columns}

  \vspace{0.05cm}
  \begin{exampleblock}{Producto de polinomios: distributiva ampliada}
    \tiny
    \[ (2x-3)(x^{2}+4x-5) = 2x^{3}+8x^{2}-10x \;-\; 3x^{2}-12x+15
       = \destacar{2x^{3}+5x^{2}-22x+15} \]
    Cada término del primero multiplica a \azul{todos} los del segundo; al
    final se reducen semejantes.
  \end{exampleblock}
\end{frame}

\begin{frame}{División de polinomios}
  \scriptsize
  \begin{columns}[T,onlytextwidth]
    \begin{column}{0.52\textwidth}
      \begin{exampleblock}{Ejemplo guiado}
        \tiny
        \[ \left(6x^{3}-x^{2}-10x+8\right) \div (2x-3) \]
        \textbf{1.} $6x^{3}\div 2x = 3x^{2}$; se resta
        $3x^{2}(2x-3)=6x^{3}-9x^{2}$ \\
        \hspace*{0.3cm} queda $8x^{2}-10x+8$

        \textbf{2.} $8x^{2}\div 2x = 4x$; se resta $8x^{2}-12x$ \\
        \hspace*{0.3cm} queda $2x+8$

        \textbf{3.} $2x\div 2x = 1$; se resta $2x-3$ \\
        \hspace*{0.3cm} queda $\destacar{11}$

        \[ \text{Cociente: } 3x^{2}+4x+1 \quad
           \text{Residuo: } 11 \]
      \end{exampleblock}
    \end{column}
    \begin{column}{0.44\textwidth}
      \begin{alertblock}{La verificación}
        \tiny
        \[ D = d\cdot c + r \]
        \[ (2x-3)(3x^{2}+4x+1) + 11 \]
        \[ = 6x^{3}-x^{2}-10x-3+11 \]
        \[ = 6x^{3}-x^{2}-10x+8 \ \checkmark \]
        \azul{Es la misma verificación de la división de enteros}, con
        polinomios.
      \end{alertblock}
    \end{column}
  \end{columns}

  \vspace{0.05cm}
  {\tiny\color{uniGris} Antes de dividir: \destacar{ordene el dividendo por
  grados descendentes} y deje un hueco (o un cero) por cada grado que falte.
  Ese es el origen de casi todos los errores.}
\end{frame}

\seccion{3.2 Productos notables}

\begin{frame}{Los cinco que hay que reconocer al instante}
  \scriptsize
  \begin{center}
  \begin{tabular}{@{}p{3.4cm}p{4.2cm}p{4.4cm}@{}}
    \toprule
    \textbf{Nombre} & \textbf{Fórmula} & \textbf{Ejemplo} \\
    \midrule
    Cuadrado de una suma &
      $(a+b)^{2} = a^{2}+2ab+b^{2}$ &
      $(x+7)^{2} = x^{2}+14x+49$ \\
    \addlinespace
    Cuadrado de una diferencia &
      $(a-b)^{2} = a^{2}-2ab+b^{2}$ &
      $(3x-4)^{2} = 9x^{2}-24x+16$ \\
    \addlinespace
    Suma por diferencia &
      $(a+b)(a-b) = a^{2}-b^{2}$ &
      $(x+5)(x-5) = x^{2}-25$ \\
    \addlinespace
    Cubo de un binomio &
      $(a+b)^{3} = a^{3}+3a^{2}b+3ab^{2}+b^{3}$ &
      $(x+2)^{3} = x^{3}+6x^{2}+12x+8$ \\
    \addlinespace
    Binomios con término común &
      $(x+m)(x+n) = x^{2}+(m+n)x+mn$ &
      $(x+3)(x+8) = x^{2}+11x+24$ \\
    \bottomrule
  \end{tabular}
  \end{center}

  \vspace{0.05cm}
  {\tiny\color{uniGris} \destacar{No son fórmulas nuevas:} son productos que
  usted ya sabe hacer con la distributiva. Se memorizan \azul{para ganar
  tiempo} y, sobre todo, porque leídas al revés son \destacar{la factorización}
  de la próxima sesión.}
\end{frame}

\begin{frame}{El término del medio: el error de todo el semestre}
  \scriptsize
  \begin{columns}[T,onlytextwidth]
    \begin{column}{0.44\textwidth}
      \centering
      \begin{tikzpicture}[scale=0.78]
        \draw[uniAzul,thick,fill=uniAzul!12] (0,0) rectangle (3,3);
        \draw[uniGris] (2,0) -- (2,3);
        \draw[uniGris] (0,2) -- (3,2);
        \node[font=\tiny] at (1,1) {$a^{2}$};
        \node[font=\tiny] at (2.5,1) {$ab$};
        \node[font=\tiny] at (1,2.5) {$ab$};
        \node[font=\tiny] at (2.5,2.5) {$b^{2}$};
        \draw[uniTitulo,<->] (0,-0.3) -- (2,-0.3)
          node[midway,below,font=\tiny,color=uniTitulo] {$a$};
        \draw[uniTitulo,<->] (2,-0.3) -- (3,-0.3)
          node[midway,below,font=\tiny,color=uniTitulo] {$b$};
      \end{tikzpicture}
    \end{column}
    \begin{column}{0.54\textwidth}
      \begin{alertblock}{Por qué $(a+b)^{2}\neq a^{2}+b^{2}$}
        \tiny
        El cuadrado de lado $a+b$ tiene área $(a+b)^2$. Al partirlo aparecen
        \destacar{cuatro} piezas: $a^2$, $b^2$ y \azul{dos rectángulos $ab$}.
        \[ (a+b)^{2} = a^{2}+\destacar{2ab}+b^{2} \]
        Los dos rectángulos del medio son el término que todo el mundo olvida.
      \end{alertblock}
      \begin{exampleblock}{Comprobación numérica}
        \tiny
        Con $a=3$, $b=4$: \ $(3+4)^2 = 49$ \ y \
        $9 + 2(12) + 16 = 49$ \checkmark \\
        Sin el término del medio daría $25$.
      \end{exampleblock}
    \end{column}
  \end{columns}
\end{frame}

\begin{frame}{Ejercicio 1 — en clase}
  \begin{alertblock}{Individual (12 minutos)}
    \scriptsize
    Desarrolle aplicando el producto notable que corresponda:

    \vspace{0.25cm}
    \centering
    \begin{tabular}{@{}lll@{}}
      1. \ $(x+9)^{2}$ &
      2. \ $(2a-5)^{2}$ &
      3. \ $(m+6)(m-6)$ \\[0.9em]
      4. \ $(2a+3b)^{2}$ &
      5. \ $(5m-2n)(5m+2n)$ &
      6. \ $(x+4)(x+7)$ \\[0.9em]
      7. \ $(x-4)(x+9)$ &
      8. \ $(x+2)^{3}$ &
      9. \ $(2x-1)^{3}$ \\
    \end{tabular}
  \end{alertblock}
\end{frame}

\begin{frame}{Ejercicio 1 — solución}
  \scriptsize
  \begin{center}
  \begin{tabular}{@{}cl@{\hspace{0.8cm}}cl@{}}
    \toprule
    \textbf{\#} & \textbf{Resultado} & \textbf{\#} & \textbf{Resultado} \\
    \midrule
    1 & $x^{2}+18x+81$ & 6 & $x^{2}+11x+28$ \\
    2 & $4a^{2}-20a+25$ & 7 & $x^{2}+5x-36$ \\
    3 & $m^{2}-36$ & 8 & $x^{3}+6x^{2}+12x+8$ \\
    4 & $4a^{2}+12ab+9b^{2}$ & 9 & $8x^{3}-12x^{2}+6x-1$ \\
    5 & $25m^{2}-4n^{2}$ & & \\
    \bottomrule
  \end{tabular}
  \end{center}

  \vspace{0.1cm}
  \begin{alertblock}{Los que hay que revisar}
    \tiny
    \textbf{4.} El doble producto es $2(2a)(3b) = 12ab$: \destacar{los
    coeficientes también entran}. \quad
    \textbf{7.} $(-4)(+9) = -36$ y $-4+9 = +5$: \azul{signos distintos}, hay
    que cuidarlos. \quad
    \textbf{9.} Los signos del cubo de una diferencia alternan:
    $+\,-\,+\,-$.
  \end{alertblock}
\end{frame}

\begin{frame}{Ejercicio 2 — en clase}
  \begin{alertblock}{En parejas (18 minutos)}
    \scriptsize
    \begin{enumerate}\setlength{\itemsep}{0.08cm}
      \item Reduzca: \ $(5x^{2}-3x+7) - (2x^{2}+8x-4)$
      \item Multiplique: \ $(3a^{2}b^{3})(-2a^{4}b)$
      \item Multiplique: \ $(x-4)(2x^{2}-3x+1)$
      \item Divida: \ $\left(x^{3}+2x^{2}-5x-6\right) \div (x+3)$ \ y
            verifique.
      \item Desarrolle y simplifique: \ $(x+3)^{2} - (x-3)^{2}$
      \item Un cuadrado tiene lado $x+5$. Escriba su área y su perímetro, y
            calcule ambos para $x = 7$.
    \end{enumerate}
  \end{alertblock}
\end{frame}

\begin{frame}{Ejercicio 2 — solución}
  \tiny
  \begin{columns}[T,onlytextwidth]
    \begin{column}{0.49\textwidth}
      \begin{block}{Puntos 1 a 3}
        \tiny
        \textbf{1.} $5x^2-3x+7-2x^2-8x+4
        = \destacar{3x^{2}-11x+11}$

        \vspace{0.08cm}
        \textbf{2.} $\destacar{-6a^{6}b^{4}}$

        \vspace{0.08cm}
        \textbf{3.} $2x^3-3x^2+x-8x^2+12x-4
        = \destacar{2x^{3}-11x^{2}+13x-4}$
      \end{block}
    \end{column}
    \begin{column}{0.49\textwidth}
      \begin{block}{Puntos 4 a 6}
        \tiny
        \textbf{4.} Cociente $\destacar{x^{2}-x-2}$, residuo $0$.
        Verificación: $(x+3)(x^2-x-2) = x^3+2x^2-5x-6$ \checkmark

        \vspace{0.08cm}
        \textbf{5.} $(x^2+6x+9)-(x^2-6x+9) = \destacar{12x}$

        \vspace{0.08cm}
        \textbf{6.} $A = (x+5)^{2} = x^2+10x+25$; \ $P = 4(x+5) = 4x+20$.
        Con $x=7$: $A = \destacar{144}$, $P = \destacar{48}$.
      \end{block}
    \end{column}
  \end{columns}

  \vspace{0.08cm}
  \begin{alertblock}{El punto 5 tiene una lectura}
    \tiny
    El resultado $12x$ \destacar{no tiene término cuadrático}: los $x^2$ se
    cancelaron. \azul{Desarrollar antes de restar} evita el error de escribir
    $(x+3)^2-(x-3)^2 = (2\cdot 3)^2$, que no significa nada.
  \end{alertblock}
\end{frame}

\begin{frame}{Ejercicio 3 — en clase (taller principal)}
  \begin{alertblock}{En grupos de 3 (20 minutos)}
    \tiny
    \setbeamerfont{itemize/enumerate body}{size=\tiny}
    \begin{enumerate}\setlength{\itemsep}{0.05cm}
      \item Desarrolle: \ $(3x+2y)^{2}$, \ $(4-x)^{2}$, \
            $(a^{2}+5)(a^{2}-5)$.
      \item Use suma por diferencia para calcular \destacar{mentalmente}
            $103 \times 97$ y $52 \times 48$.
      \item Escriba como producto notable y desarrolle: el área de un cuadrado
            cuyo lado mide $2x-7$.
      \item Un terreno rectangular mide $(x+4)$ por $(x+9)$ metros. Escriba el
            área como polinomio y calcule para $x=11$.
      \item Encuentre el error: \ $(2x+3)^{2} = 4x^{2}+9$
      \item Encuentre el error: \ $(x-5)(x+5) = x^{2}+25$
    \end{enumerate}
  \end{alertblock}
\end{frame}

\begin{frame}{Ejercicio 3 — solución}
  \scriptsize
  \begin{columns}[T,onlytextwidth]
    \begin{column}{0.49\textwidth}
      \begin{block}{Puntos 1 a 4}
        \tiny
        \textbf{1.} $9x^2+12xy+4y^2$; \ $16-8x+x^2$; \ $a^4-25$

        \vspace{0.06cm}
        \textbf{2.} $103\cdot 97 = (100+3)(100-3) = 10\,000-9
        = \destacar{9\,991}$ \\
        $52\cdot 48 = (50+2)(50-2) = 2\,500-4 = \destacar{2\,496}$

        \vspace{0.06cm}
        \textbf{3.} $(2x-7)^{2} = \destacar{4x^{2}-28x+49}$

        \vspace{0.06cm}
        \textbf{4.} $A = x^2+13x+36$; con $x=11$: $\destacar{300}$ m$^2$
      \end{block}
    \end{column}
    \begin{column}{0.49\textwidth}
      \begin{alertblock}{Puntos 5 y 6}
        \tiny
        \textbf{5.} Falta el doble producto:
        $(2x+3)^{2} = 4x^{2}+\destacar{12x}+9$.

        \vspace{0.08cm}
        \textbf{6.} El signo:
        $(x-5)(x+5) = x^{2}-\destacar{25}$. En suma por diferencia el resultado
        \azul{siempre} lleva menos.
      \end{alertblock}
    \end{column}
  \end{columns}

  \vspace{0.05cm}
  {\tiny\color{uniGris} El punto \destacar{2} muestra para qué sirve de verdad
  un producto notable: convierte una multiplicación difícil en una resta
  trivial. \azul{Es la misma idea de los logaritmos: bajar el nivel de la
  operación.}}
\end{frame}

\begin{frame}{Errores frecuentes en esta sesión}
  \scriptsize
  \begin{center}
  \begin{tabular}{@{}p{5.2cm}p{6.6cm}@{}}
    \toprule
    \textbf{Error} & \textbf{Corrección} \\
    \midrule
    $(a+b)^{2} = a^{2}+b^{2}$ &
      Falta $2ab$: son \azul{tres} términos \\
    \addlinespace
    $(a-b)(a+b) = a^{2}+b^{2}$ &
      Es $a^{2}-b^{2}$ \\
    \addlinespace
    $(2x)^{2} = 2x^{2}$ &
      Es $4x^{2}$: el coeficiente también se eleva \\
    \addlinespace
    Restar sin cambiar todos los signos &
      $-(x^{2}-7x) = -x^{2}+7x$ \\
    \addlinespace
    Sumar $3x^{2}+5x$ como $8x^{3}$ &
      No son semejantes: se deja indicado \\
    \addlinespace
    Dividir sin ordenar el dividendo &
      Se ordena por grados y se completan los que falten \\
    \bottomrule
  \end{tabular}
  \end{center}
\end{frame}

\begin{frame}{Cierre de la sesión}
  \begin{block}{Lo que debe quedar claro hoy}
    \scriptsize
    \begin{itemize}\setlength{\itemsep}{0.02cm}
      \item Solo se suman \azul{términos semejantes}; multiplicar siempre se
            puede.
      \item Producto de polinomios: distributiva aplicada término a término.
      \item División: se ordena, se divide, se resta y se verifica con
            $D = d\cdot c + r$.
      \item Los cinco productos notables, con \destacar{el término del medio}
            incluido.
      \item Suma por diferencia da $a^{2}-b^{2}$: \azul{siempre con menos}.
      \item Leídos al revés, los productos notables son la factorización.
    \end{itemize}
  \end{block}

  \vspace{0.08cm}
  \begin{alertblock}{Trabajo independiente}
    \tiny
    Resolver los 20 ejercicios de la lámina siguiente. En la sesión 14 vemos
    \destacar{factorización}: exactamente estos mismos productos, recorridos en
    sentido contrario.
  \end{alertblock}
\end{frame}

\begin{frame}{Ejercicios propuestos}
  \vspace{-0.1cm}
  \begin{beamercolorbox}[wd=\textwidth,sep=0.10cm]{block body}
    \centering\scriptsize
    \textbf{1--7:} operaciones. \quad \textbf{8--15:} productos notables.
    \quad \textbf{16--20:} aplicaciones.
  \end{beamercolorbox}
  \listatiny
  \vspace{0.15cm}
  \begin{columns}[T,onlytextwidth]
    \begin{column}{0.47\textwidth}
      \begin{enumerate}\setlength{\itemsep}{0.16cm}
        \item $(2x^{2}+5x-1)+(3x^{2}-2x+6)$
        \item $(7a-4)-(2a-9)$
        \item $(5x^{3}y)(-3xy^{2})$
        \item $2x(3x^{2}-4x+1)$
        \item $(x+2)(x^{2}-3x+5)$
        \item $(x^{2}+3x-10)\div(x-2)$
        \item Grado de $-4x^{3}y^{2}z$
        \item $(x+11)^{2}$
        \item $(3m-2)^{2}$
        \item $(y+8)(y-8)$
      \end{enumerate}
    \end{column}
    \begin{column}{0.50\textwidth}
      \begin{enumerate}\setlength{\itemsep}{0.16cm}
        \setcounter{enumi}{10}
        \item $(4a+5b)^{2}$
        \item $(x-1)^{3}$
        \item $(x+6)(x+2)$
        \item $(x-7)(x+3)$
        \item $(2x^{2}+3)(2x^{2}-3)$
        \item Calcule $98\times 102$ con un producto notable
        \item Calcule $41^{2}$ como $(40+1)^{2}$
        \item Área de un cuadrado de lado $3x-1$
        \item Área de un rectángulo de lados $(x+5)$ y $(x-5)$
        \item ¿Por qué $(a+b)^{2}\neq a^{2}+b^{2}$?
      \end{enumerate}
    \end{column}
  \end{columns}
\end{frame}

\begin{frame}{Respuestas}{Autoevaluación de los ejercicios propuestos}
  \listatiny
  \vspace{0.15cm}
  \begin{columns}[T,onlytextwidth]
    \begin{column}{0.47\textwidth}
      \begin{enumerate}\setlength{\itemsep}{0.16cm}
        \item $5x^{2}+3x+5$
        \item $5a+5$
        \item $-15x^{4}y^{3}$
        \item $6x^{3}-8x^{2}+2x$
        \item $x^{3}-x^{2}-x+10$
        \item $x+5$, residuo $0$
        \item $6$
        \item $x^{2}+22x+121$
        \item $9m^{2}-12m+4$
        \item $y^{2}-64$
      \end{enumerate}
    \end{column}
    \begin{column}{0.50\textwidth}
      \begin{enumerate}\setlength{\itemsep}{0.16cm}
        \setcounter{enumi}{10}
        \item $16a^{2}+40ab+25b^{2}$
        \item $x^{3}-3x^{2}+3x-1$
        \item $x^{2}+8x+12$
        \item $x^{2}-4x-21$
        \item $4x^{4}-9$
        \item $(100-2)(100+2) = 9\,996$
        \item $1\,600+80+1 = 1\,681$
        \item $9x^{2}-6x+1$
        \item $x^{2}-25$
        \item Porque falta el doble producto $2ab$
      \end{enumerate}
    \end{column}
  \end{columns}
\end{frame}

\begin{frame}{Referencias}
  \footnotesize
  \begin{itemize}
    \item[{[1]}] Ramos del Olmo, Á. M. (2016). \textit{Matemáticas básicas
      para el acceso a la universidad}. Ediciones Pirámide.
      \texttt{https://elibro.net/es/ereader/uniminuto/49134}
    \item[{[2]}] González Pulido, J. W., Chacón Benavides, J. A. \&
      Fonseca Correa, L. Á. (2023). \textit{Matemática básica} (1.\textsuperscript{a} ed.).
      Universidad Mariana.
      \texttt{https://elibro.net/es/ereader/uniminuto/280647}
    \item[{[3]}] Placencia Valero, J. (2008). \textit{Compendio de matemática
      básica elemental} (1.\textsuperscript{a} ed.). Editorial Tébar Flores.
      \texttt{https://elibro.net/es/ereader/uniminuto/269758}
  \end{itemize}

  \vspace{0.35cm}
  \begin{center}
    {\Large\color{uniAzul}\textbf{¿Preguntas?}}
  \end{center}
\end{frame}

\end{document}
